\documentclass[10pt,a4paper]{amsart}
\usepackage[margin=19mm]{geometry}
\usepackage{amsmath,amssymb,mathtools}
\usepackage[scr=rsfs,cal=euler]{mathalfa}
\usepackage{xfrac}
\usepackage[unicode,hidelinks,bookmarks=false]{hyperref}
\newcommand{\ie}{i.e.\ }
\newcommand{\cf}{cf.\ }
\newcommand{\resp}{respectively\ }
\newcommand{\Subsection}{\subsection}
\providecommand{\nobreakdash}{\nobreak\hbox{-}}
\setlength{\emergencystretch}{2em}
\multlinegap0pt
\usepackage{bm}
\usepackage{amssymb}
\usepackage{mathtools}
\usepackage{enumerate}
\usepackage{dsfont}
\newtheorem{lemma}{Lemma}
\newcommand{\Ordo}{\mathcal{O}}
\renewcommand{\d}{\,\mathrm{d}}
\renewcommand{\i}{\mathrm{i}}
\DeclareMathOperator{\one}{\mathds{1}}
\title{Marked GUE-corners process in doubly periodic dimer models}
\author{Tomas Berggren, Nedialko Bradinoff}
\date{Source: arXiv:2603.27906v1 (CC BY 4.0)}
\begin{document}
\maketitle

\noindent\emph{Source and licence.} Marked GUE-corners process in doubly periodic dimer models. Version arXiv:2603.27906v1 (CC BY 4.0), distributed by arXiv under the Creative Commons Attribution 4.0 International licence, http://creativecommons.org/licenses/by/4.0/, as recorded in the arXiv licence metadata for this exact version and shown on the abstract page https://arxiv.org/abs/2603.27906v1. Full licence notice: \texttt{licenses/arxiv-2603.27906-CC-BY-4.0.txt}.\par\smallskip

\noindent\textit{Editorial scope.} Counted unit: the article's lemma lemma:Js (source0.tex lines 2434--2467). The article's complete original proof of it is delivered with it (source0.tex lines 2470--2562, one \texttt{proof} environment of 93 lines). No mathematics is written, repaired or re-derived: the statement and the proof are the article's own lines, in the article's own notation, and no retained line is substituted or re-flowed.  Retained uncounted as the source-local context the proof needs: lines 984--990; lines 1573--1604; lines 1661--1674; lines 1815--1831; lines 1834--1854; lines 1927--1932.  Nothing was omitted from any retained span, so the package carries no omission record.  No retained line contains a citation, so the wrapper carries no bibliography and no bibliography entry.  No retained line contains a figure environment, an image-inclusion command or drawing code, so no image is delivered and the package declares no asset dependency.  Editorial formatting only, in addition: an AMS \texttt{amsart} preamble that restates the article's own declarations and macros used by the retained lines, namely the environments \texttt{lemma} on separate counters, together with the standard packages those lines need.  The retained environments print in this document as Lemma 1 in order of appearance, whereas the article prints the same items with the numbering of its own full document.  The complete native source of the article as distributed in the arXiv e-print for this exact version is bundled unchanged as \texttt{sources/197397/source0.tex}.
\begin{equation}\label{eq:period_function}
\nu(t,j)=
\begin{cases}
\frac{1}{1+\alpha_{\ell+1-t}\beta_{\ell-t}^{-1}}, & j=0, \\
\frac{\alpha_{\ell+1-t}\beta_{\ell-t}^{-1}}{1+\alpha_{\ell+1-t}\beta_{\ell-t}^{-1}}, & j=1.
\end{cases}
\end{equation}

\begin{lemma}\label{lem:at_1}
The following identities hold:
\begin{equation}\label{eq:at_1_1}
Q(q_\infty)=\frac{1}{1+\beta_\ell^{-1} \alpha_1}
\begin{pmatrix}
1 \\
\alpha_1
\end{pmatrix}
\begin{pmatrix}
1 & \beta_\ell^{-1}
\end{pmatrix},
\end{equation}
\begin{equation}\label{eq:at_1_2}
\left.(z-1)^\ell w\right|_{(z,w)=q_{\infty}}=\prod_{m=1}^\ell(1+\alpha_m^{-1}\beta_m)(1+\alpha_{m+1}\beta_m^{-1}),
\end{equation}
\begin{equation}\label{eq:at_1_2_q0}
\left.(z-1)^{-\ell} w\right|_{(z,w)=q_{0}}=\left(\prod_{m=1}^\ell(1+\alpha_m^{-1}\beta_m)(1+\alpha_{m+1}\beta_m^{-1})\right)^{-1},
\end{equation}
and
\begin{equation}\label{eq:at_1_3}
\left.(z-1)^{i-i'}\prod_{m=2i'+1}^{2i}\phi_m(z)\right|_{z=1}
=\prod_{m=i'+1}^i(1+\alpha_m^{-1}\beta_m)\prod_{m=i'+1}^{i-1}(1+\alpha_{m+1}\beta_m^{-1})
\begin{pmatrix}
1 \\
\alpha_{i'+1}
\end{pmatrix}
\begin{pmatrix}
1 & \beta_i^{-1}
\end{pmatrix},
\end{equation}
for~$0\leq i'<i\leq \ell$.
\end{lemma}

\begin{lemma}\label{lem:matrix_local}
For $(z_1,w_1)$ and $(z_2,w_2)$ in a neighbourhood of $q_{\infty}$, let
~$z_i=1+\frac{Z_i}{N^\frac{1}{2}}$,~$i=1,2$, and set
\begin{equation}
\tilde g(\ell x+i,2y+j)=\prod_{m=1}^i(1+\alpha_m^{-1}\beta_m)(1+\alpha_{m+1}\beta_m^{-1})\alpha_{i+1}^{-j},
\end{equation}
and let~$\nu$ be given in~\eqref{eq:period_function}. Then, as~$N\to \infty$, 
\begin{multline} \label{eqn:matrix_precise}
\boldsymbol{\mathcal{M}}_{i_1,j_1,i_2,j_2}\left(z_1,w_1;z_2,w_2\right)
%
= \nu(\ell-i_2,j_2)\frac{Z_1^{i_1}}{Z_2^{i_2}}\frac{\tilde g(\ell x_2+i_2,2y_2+j_2)}{\tilde g(\ell x_1+i_1,2y_1+j_1)}\frac{N^{\frac{i_2}{2}}}{N^{\frac{i_1}{2}}}\left(1+\Ordo\left(\frac{|Z_1|+|Z_2|}{N^{\frac{1}{2}}}\right)\right),
\end{multline}
where, assuming $|Z_k|\leq N^{\frac{1}{2}}$, the implicit constant in the remainder term is independent of $Z_k, x_k, i_k, y_k, j_k$, $k=1,2$.
\end{lemma}

\begin{lemma}\label{lemma:local_action_function}
Let~$z_k=1+\frac{Z_k}{N^\frac{1}{2}}$,~$k=1,2$. Then, as~$N\to \infty$, for $(z_1,w_1)$,  $(z_2,w_2)$ in the neighbourhood of $q_{\infty}$, we have that
\begin{multline} \label{eqn:integrand_near_q_infinity}
\exp\left({N(G(z_2,w_2)-G(z_1,w_1))}\right)
\frac{w_1^{r_1}}{w_2^{r_2}}
 \frac{z_1^{\mu_1\sqrt{N}}}{z_2^{\mu_2\sqrt{N}}}
 \frac{1}{z_2(z_2-z_1)}
\\= N^{\frac{1}{2}} \frac{\left(1+\mathcal{O}\left( \frac{|Z_1|^3+|Z_2|^3}{N^{\frac{1}{2}}} \right)\right)}{Z_2-Z_1} 
 N^{\frac{\ell r_1}{2}-\frac{\ell r_2}{2}} 
B^{ r_1- r_2} 
\frac{Z_2^{\ell r_2}}{Z_1^{\ell r_1}}
\exp\left(\frac{\sigma^2}{2}(Z_2^2-Z_1^2)+Z_1\mu_1-Z_2\mu_2\right),
\end{multline}
where $B=\left(\prod_{m=1}^\ell(1+\alpha_m^{-1}\beta_m)(1+\alpha_{m+1}\beta_m^{-1})\right)$ and the implicit constant in the remainder term is uniform for compact subsets of $\Lambda_{\{0,1\}}^2$. 
 For $|Z_1|, |Z_2| \leq N^{\frac{1}{12}}$, the remainder term
 is uniformly $\mathcal{O}(N^{-\frac{1}{4}})$, as $N\rightarrow\infty$.
\end{lemma}

\begin{lemma}\label{lemma:expansions}
Let~$z_k=1+\frac{Z_k}{N^\frac{1}{2}}$,~$k=1,2$, with $|Z_k|< N^{\frac{1}{2}}$ (on the canonical local chart in this text, at a neighbourhood of $q_{\infty}$). Then, as~$N\to \infty$, we have that
\begin{equation} \label{eqn:local_action_contribution}
\exp(NG(z_k,w_k))= \exp\left(NG(1)+\frac{\sigma^2}{2}Z_k^2+\mathcal{O}\left(\frac{|Z_k|^3}{N^{\frac{1}{2}}}\right)\right),
\end{equation}
\begin{equation} \label{eqn:local_w_contribution}
w_k^{r_k}=
Z_k^{-\ell r_k} 
N^{\frac{\ell r_k}{2}} 
\left(B+\mathcal{O}\left(\frac{|Z_k|}{N^{\frac{1}{2}}}\right)\right)^{r_k}
=
Z_k^{-\ell r_k} 
N^{\frac{\ell r_k}{2}} 
B^{r_k}\left(1+\mathcal{O}\left(\frac{r_k |Z_k|}{N^{\frac{1}{2}}}\right)\right),
\end{equation}
and
\begin{equation} \label{eqn:local_z_contribution}
z_i^{\sqrt{N}\mu_k}=\left(1+\frac{Z_k}{ \sqrt{N}}\right)^{\mu_k\sqrt{N}}=\exp\left(\mu_k Z_k+\mathcal{O}\left(\frac{\mu_k |Z_k|^2}{N^{\frac{1}{2}}}\right)\right),
\end{equation}
where $B$ is as in Lemma \ref{lemma:local_action_function}.
\end{lemma}

\begin{equation}\label{eqn:gauge_proof_coordinates}
\frac{h(\ell x_1+i_1,2y_1+j_1)}{h(\ell x_2+i_2,2y_2+j_2)}
=N^{\frac{\ell r_2-i_2}{2}-\frac{\ell r_1-i_1}{2}} 
B^{r_2-r_1}
\frac{\tilde g(\ell x_1+i_1,2y_1+j_1)}{\tilde g(\ell x_2+i_2,2y_2+j_2)}.
\end{equation}

\begin{lemma}%[Computing the single integral] 
\label{lemma:Js}
Assuming $\ell r_1-i_1> \ell r_2-i_2$, we have that 
\begin{multline}
N^{\frac{1}{2}}
\frac{h(\ell x_1+i_1,2y_1+j_1)}{h(\ell x_2+i_2,2y_2+j_2)}
 \mathcal{J}_{s}=
%
%
%
\\=\begin{cases}
\mathcal{O}
\left(N^{\frac{\ell r_2-i_2}{2}-\frac{\ell r_1-i_1}{2}+\frac{1}{2}} \exp\left(-\sqrt{N}\log 2(\mu_1-\mu_2)\right)\right) 
 &\text{for } \mu_1 \leq \mu_2,
%\\
%\quad\text{for}\quad \mu_1=\mu_2
\\
\frac{\nu(\ell-i_2,j_2)}{2\pi i}
%
%
\int_{|Z|=1} 
  Z^{\ell r_2-i_2-\ell r_1+i_1}
\exp((\mu_1-\mu_2) Z)
dZ
+
%
\mathcal{O}\left(N^{-\frac{1}{2}}\right)
&\text{for} \quad \mu_1>\mu_2,
\end{cases}
\end{multline}
as $N\rightarrow\infty$ and in particular the left-hand side is uniformly bounded in $N$ on compact subsets of $\Lambda_{\{0,1\}}^2$, under the constructed embedding, and converges point-wise.
%
%
\end{lemma}

\begin{proof}
From Lemmas \ref{lem:at_1} and \ref{lem:matrix_local} we can deduce that
$\boldsymbol{\mathcal{M}}_{i_1,j_1,i_2,j_2}(z,w;z,w) 
w^{r_1-r_2}$
has a pole of order $(\ell r_1-i_1)-(\ell r_2-i_2)$ at $q_{\infty}$ and a zero at $q_0$.
\begin{enumerate}
\item
Now if $\mu_1 \leq \mu_2$, since $q_{\infty}$ and $p_0$ are in the interior of the curve $\tilde{\Gamma}_s$ and the integrand is analytic at $q_0$, we deform the contour so that $|z|=2$. Hence we have that 
\begin{equation}
(2\pi)^{-1}\left|\boldsymbol{\mathcal{M}}_{i_1,j_1,i_2,j_2}(z,w;z,w) 
w^{r_1-r_2}
z^{(\mu_1-\mu_2)\sqrt{N}-1}\right|
\leq H(\ell r_1, i_1, \ell r_2, i_2) 2^{-\sqrt{N}(\mu_2-\mu_1)-1}
,
\end{equation}
where $H$ is uniformly bounded on compact sets of $\Lambda_{\{0,1\}}$. 
Hence since the length of the curve of integration is finite and recalling the form of the gauge factor, \eqref{eqn:gauge_proof_coordinates},  
there is a constant $C>0$ (independent of $N$) such that
$$N^{\frac{1}{2}}
\frac{h(\ell x_1+i_1,2y_1+j_1)}{h(\ell x_2+i_2,2y_2+j_2)}  |\mathcal{J}_s|\leq 
 C B^{r_2-r_1} 
N^{\frac{\ell r_2-i_2}{2}-\frac{\ell r_1-i_1}{2}+\frac{1}{2}}
 H(\ell r_1, i_1, \ell r_2, i_2)  2^{-\sqrt{N}(\mu_2-\mu_1)-1}.
 $$
In particular since $\ell r_1-i_1-\ell r_2 +i_2 \geq 1$ we have a uniform bound on $\mathcal{J}_s$ for $\mu_1\leq \mu_2$ (on compact subsets of $\Lambda_{\{0,1\}}^2$). %
%
Pointwise, for $\mu_1<\mu_2$, we have that
\begin{equation}
N^{\frac{1}{2}}
\frac{h(\ell x_1+i_1,2y_1+j_1)}{h(\ell x_2+i_2,2y_2+j_2)}
 |\mathcal{J}_{s}|
=\mathcal{O}
\left(\exp\left(-\sqrt{N}\log 2(\mu_2-\mu_1)\right)\right), \quad \text{as}\quad N\rightarrow\infty.
\end{equation}
\item
Let $\mu_1\geq \mu_2$. In this case the pole at $p_0$ is finite and we can deform $\Gamma_s$ to a simple closed curve treading around $p_0$, $\Gamma_s^{(1)}$ with $p_{\infty}, q_{0}, q_{\infty}$ in its exterior, and a simple curve around $q_{\infty}$, $\Gamma_s^{(2)}$ with $p_0$, $q_0$, $p_{\infty}$ in its exterior. 

%
%
 We choose  $\Gamma_s^{(2)}$ to be a circle of radius $\frac{1}{\sqrt{N}}$ around $q_{\infty}$, and we make the change of variables $z=1+\frac{Z}{\sqrt{N}}, dz=\frac{dZ}{\sqrt{N}}$ (on the circle $|Z|=1$). By the identities proved in Lemmas \ref{lem:matrix_local} and \ref{lemma:expansions}, we have that
 %
\begin{multline}
N^{\frac{1}{2}}
\frac{h(\ell x_1+i_1,2y_1+j_1)}{h(\ell x_2+i_2,2y_2+j_2)}
%\mathcal{J}_s
\frac{\one_{\ell r_1-i_1>\ell r_2-i_2}}{(2\pi\i)}\int_{ \Gamma_s^{(2)}}
\boldsymbol{\mathcal{M}}_{i_1,j_1,i_2,j_2}(z,w;z,w) 
w^{r_1-r_2}
z^{(\mu_1-\mu_2)\sqrt{N}}
\frac{\d z}{z}
%
%
\\=
%\left( N^{\frac{i_2}{2}-\frac{i_1}{2}} Z^{i_1-i_2} 
%
\frac{\nu(\ell-i_2,j_2)}{2\pi i}
%
\int_{|Z|=1} Z^{\ell r_2-i_2-\ell r_1+i_1}
\exp((\mu_1-\mu_2) Z)
 \left(1+\mathcal{O}(N^{-\frac{1}{2}})\right)dZ
% \\=\boldsymbol{M} B^{r_1-r_2} N^{\frac{\ell r_2-i_2-\ell r_1+i_1}{2}} N^{-1/2}
\\= \nu(\ell-i_2,j_2)
%
\left(
 \frac{1}{2\pi i}\int_{|Z|=1} 
  Z^{\ell r_2-i_2-\ell r_1+i_1}
\exp((\mu_1-\mu_2) Z)
\d Z
+\mathcal{O}\left(N^{-\frac{1}{2}}\right)\right),
\text{  as  }
 N\rightarrow\infty,
  \end{multline}
where the implicit constant is uniformly bounded on compact subsets of $\Lambda_{\{0,1\}}^2$.


The integral over $\Gamma_s^{(1)}$ is uniformly bounded, since the pole at $p_0$ of the integrand is of finite order (in fact $\leq 1$). Thus we have that uniformly on compact subsets of $\Lambda_{\{0,1\}}^2$,
\begin{multline}\label{eqn:justify_last_line}
N^{\frac{1}{2}}
\frac{h(\ell x_1+i_1,2y_1+j_1)}{h(\ell x_2+i_2,2y_2+j_2)}
\frac{\one_{\ell r_1-i_1>\ell r_2-i_2}}{(2\pi\i)}\int_{ \Gamma_s^{(1)}}
\boldsymbol{\mathcal{M}}_{i_1,j_1,i_2,j_2}(z,w;z,w) 
w^{r_1-r_2}
z^{(\mu_1-\mu_2)\sqrt{N}}
\frac{\d z}{z}
\\=\mathcal{O}\left(N^{\frac{\ell r_2-i_2-\ell r_1+i_1+1}{2}}\right),\quad\text{as } N\rightarrow\infty.
\end{multline}
%
%
If $\mu_1>\mu_2$, for $N$ sufficiently large, the integrand of $\mathcal{J}_s$ is analytic at $p_0$ and so the integral over $\Gamma_s^{(1)}$ is 0.
\end{enumerate}
The proof is now complete.
% 
\end{proof}

\end{document}
